%%%%%%%%%%%%%%%%%
% This is an sample CV template created using altacv.cls
% (v1.7.4, 30 July 2025) written by LianTze Lim (liantze@gmail.com). Compiles with pdfLaTeX, XeLaTeX and LuaLaTeX.
%
%% It may be distributed and/or modified under the
%% conditions of the LaTeX Project Public License, either version 1.3
%% of this license or (at your option) any later version.
%% The latest version of this license is in
%%    http://www.latex-project.org/lppl.txt
%% and version 1.3 or later is part of all distributions of LaTeX
%% version 2003/12/01 or later.
%%%%%%%%%%%%%%%%

%% Use the "normalphoto" option if you want a normal photo instead of cropped to a circle
% \documentclass[10pt,a4paper,withhyper,normalphoto]{altacv}

\documentclass[10pt,a4paper,withhyper]{altacv}
%% AltaCV uses the fontawesome5 and simpleicons packages.
%% See http://texdoc.net/pkg/fontawesome5 and http://texdoc.net/pkg/simpleicons for full list of symbols.

% Change the page layout if you need to
\geometry{left=1.25cm,right=1.25cm,top=1.5cm,bottom=1.5cm,columnsep=1.2cm}

% The paracol package lets you typeset columns of text in parallel
\usepackage{paracol}

% Change the font if you want to, depending on whether
% you're using pdflatex or xelatex/lualatex
% WHEN COMPILING WITH XELATEX PLEASE USE
% xelatex -shell-escape -output-driver="xdvipdfmx -z 0" sample.tex
\iftutex
  % If using xelatex or lualatex:
  \setmainfont{Roboto Slab}
  \setsansfont{Lato}
  \renewcommand{\familydefault}{\sfdefault}
\else
  % If using pdflatex:
  \usepackage[rm]{roboto}
  \usepackage[defaultsans]{lato}
  % \usepackage{sourcesanspro}
  \renewcommand{\familydefault}{\sfdefault}
\fi

% Change the colours if you want to
\definecolor{SlateGrey}{HTML}{2E2E2E}
\definecolor{LightGrey}{HTML}{666666}
\definecolor{DarkPastelRed}{HTML}{450808}
\definecolor{PastelRed}{HTML}{8F0D0D}
\definecolor{GoldenEarth}{HTML}{E7D192}
\colorlet{name}{black}
\colorlet{tagline}{PastelRed}
\colorlet{heading}{DarkPastelRed}
\colorlet{headingrule}{GoldenEarth}
\colorlet{subheading}{PastelRed}
\colorlet{accent}{PastelRed}
\colorlet{emphasis}{SlateGrey}
\colorlet{body}{LightGrey}

% Change some fonts, if necessary
\renewcommand{\namefont}{\Huge\rmfamily\bfseries}
\renewcommand{\personalinfofont}{\footnotesize}
\renewcommand{\cvsectionfont}{\LARGE\rmfamily\bfseries}
\renewcommand{\cvsubsectionfont}{\large\bfseries}


% Change the bullets for itemize and rating marker
% for \cvskill if you want to
\renewcommand{\cvItemMarker}{{\small\textbullet}}
\renewcommand{\cvRatingMarker}{\faCircle}
% ...and the markers for the date/location for \cvevent
% \renewcommand{\cvDateMarker}{\faCalendar*[regular]}
% \renewcommand{\cvLocationMarker}{\faMapMarker*}


% If your CV/résumé is in a language other than English,
% then you probably want to change these so that when you
% copy-paste from the PDF or run pdftotext, the location
% and date marker icons for \cvevent will paste as correct
% translations. For example Spanish:
% \renewcommand{\locationname}{Ubicación}
% \renewcommand{\datename}{Fecha}


%% Use (and optionally edit if necessary) this .tex if you
%% want to use an author-year reference style like APA(6)
%% for your publication list
% \input{pubs-authoryear.cfg}

%% Use (and optionally edit if necessary) this .tex if you
%% want an originally numerical reference style like IEEE
%% for your publication list
\input{pubs-num.cfg}

%% sample.bib contains your publications
\addbibresource{sample.bib}

\begin{document}
\name{Maxime Lalisse}
\tagline{}
%\tagline{Fullstack Developer}
%% You can add multiple photos on the left or right
\photoR{2.8cm}{photo}
% \photoL{2.5cm}{Yacht_High,Suitcase_High}

\personalinfo{%
  %\phone{+33 (0) 6 36 71 54 50}
  \email{maxime@lalis.se}
  \homepage{https://maxime.lalis.se}
  %\mailaddress{Åddrésş, Street, 00000 Cóuntry}
  \github{mjal}
  \location{Lille, FRANCE}
  % \twitter{@twitterhandle}
  %\xtwitter{@x-handle}

  %\linkedin{maxime-lalisse-42078a1a8}

  %\orcid{0000-0000-0000-0000}
  %% You can add your own arbitrary detail with
  %% \printinfo{symbol}{detail}[optional hyperlink prefix]
  % \printinfo{\faPaw}{Hey ho!}[https://example.com/]

  %% Or you can declare your own field with
  %% \NewInfoFiled{fieldname}{symbol}[optional hyperlink prefix] and use it:
  %\NewInfoField{gitlab}{\faGitlab}[https://gitlab.com/]
  %\gitlab{your_id}
  %%
  %% For services and platforms like Mastodon where there isn't a
  %% straightforward relation between the user ID/nickname and the hyperlink,
  %% you can use \printinfo directly e.g.
  % \printinfo{\faMastodon}{@username@instace}[https://instance.url/@username]
  %% But if you absolutely want to create new dedicated info fields for
  %% such platforms, then use \NewInfoField* with a star:
  %\NewInfoField*{mastodon}{\faMastodon}
  %% then you can use \mastodon, with TWO arguments where the 2nd argument is
  %% the full hyperlink.
  %\mastodon{@username@instance}{https://instance.url/@username}
}

\makecvheader
%% Depending on your tastes, you may want to make fonts of itemize environments slightly smaller
% \AtBeginEnvironment{itemize}{\small}

%% Set the left/right column width ratio to 6:4.
\columnratio{0.6}

% Start a 2-column paracol. Both the left and right columns will automatically
% break across pages if things get too long.
\begin{paracol}{2}
\cvsection{Experience}

\cvevent{Electronic Voting Engineer}{Gedivote}{November 2025 -- Present}{Remote}
\begin{itemize}
\item Platform development (general, security and cryptography)
\end{itemize}

\divider

\cvevent{Security Research on Internet Voting}{CNRS}{March 2025 -- August 2025}{Nancy}
\begin{itemize}
\item Design of voting protocols with ballot stuffing
\item Development of zero-knowledge proofs
\item Formal verification with ProVerif
\end{itemize}

\divider

\cvevent{Fullstack Developer - React + e-voting}{Scrutin.app}{February 2023 -- February 2025}{Remote}
\begin{itemize}
\item Design of an open-source web and mobile internet voting software
\item Development of a Belenios implementation
\end{itemize}

\divider

\cvevent{Fullstack Developer - Django + e-voting}{Electis}{July 2022 -- January 2023}{Paris}
\begin{itemize}
\item ElectionGuard integration
\item Adaptations for works council elections
\end{itemize}

\divider

\cvevent{Fullstack Developer - Ruby on rails + Vue}{Winddle}{August 2021 -- May 2022}{Paris}
Replacement of Backbone.js logic with Vue.js and SQL query optimization.

\divider

\cvevent{Cybersecurity Engineer}{Stormshield}{November 2020 -- July 2021}{Lille}
Development around TPM, IPSEC, and FreeBSD.

\divider

\cvevent{C++ Developer - Video Games}{Eko Software}{September 2014 -- October 2020}{Paris}
Development of a game engine running on multiple platforms (PS3, PS4, Xbox One, Switch). GPU and shader development. Creation of web tools for managing player databases.

\newpage

\cvsection{Projects}

%\cvevent{Project 1}{Funding agency/institution}{}{}
\cvevent{\href{https://scrutin.app}{Scrutin.app}}{}{2023 - Present}{}

An \textit{easy-to-use} and \textit{secure} internet voting platform based on the Belenios voting protocol (React Native + Node.js).
\href{https://github.com/mjal/scrutin}{(source)}
\href{https://scrutin.app/}{(demo)}
\href{https://docs.scrutin.app/}{(doc)}

\divider

\cvevent{\href{https://github.com/mjal/sirona}{Sirona}}{}{2024 - Present}{}

A tool to interact with elections created with Belenios.

\divider

\cvevent{Posca (contributor)}{}{2024}{}

A Matrix client with social network features

\divider

\cvevent{Kazarma (contributor)}{}{2023}{}

A Matrix <-> ActivityPub bridge

\divider

\cvevent{\href{https://github.com/safebook/zero}{Zero}}{}{2023}{}

A prototype of metadata-free messaging inspired by Bitmessage.

\divider

\cvevent{Ami}{}{2013 - ...}{}

A hackable homepage, self-hosted or local. Currently delegates to other sites (Wikipedia, ...) based on bangs (!w, ...). Usable as a default search engine in Firefox.
\href{https://mjal.github.io/ami/}{(v1)}
\href{https://mjal.github.io/ami3d/}{(v2)}
\href{https://mjal.github.io/ami3/}{(v3)}

\divider

\cvevent{Safebook}{}{2011 - ... (paused)}{}

An experiment in end-to-end encrypted social networking. Originally a browser extension adding an "Encrypt and send" button to Facebook and decrypting friends' messages on the fly.

\medskip

%\cvsection{A Day of My Life}
%
%% Adapted from @Jake's answer from http://tex.stackexchange.com/a/82729/226
%% \wheelchart{outer radius}{inner radius}{
%% comma-separated list of value/text width/color/detail}
%\wheelchart{1.5cm}{0.5cm}{%
%  6/8em/accent!30/{Sleep,\\beautiful sleep},
%  3/8em/accent!40/Hopeful novelist by night,
%  8/8em/accent!60/Daytime job,
%  2/10em/accent/Sports and relaxation,
%  5/6em/accent!20/Spending time with family
%I}
%
%% use ONLY \newpage if you want to force a page break for
%% ONLY the current column
%\newpage
%
%\cvsection{Publications}
%
%%% Specify your last name(s) and first name(s) as given in the .bib to automatically bold your own name in the publications list.
%%% One caveat: You need to write \bibnamedelima where there's a space in your name for this to work properly; or write \bibnamedelimi if you use initials in the .bib
%%% You can specify multiple names, especially if you have changed your name or if you need to highlight multiple authors.
%\mynames{Lim/Lian\bibnamedelima Tze,
%  Wong/Lian\bibnamedelima Tze,
%  Lim/Tracy,
%  Lim/L.\bibnamedelimi T.}
%%% MAKE SURE THERE IS NO SPACE AFTER THE FINAL NAME IN YOUR \mynames LIST
%
%\nocite{*}
%
%\printbibliography[heading=pubtype,title={\printinfo{\faBook}{Books}},type=book]
%
%\divider
%
%\printbibliography[heading=pubtype,title={\printinfo{\faFile*[regular]}{Journal Articles}},type=article]
%
%\divider
%
%\printbibliography[heading=pubtype,title={\printinfo{\faUsers}{Conference Proceedings}},type=inproceedings]

%% Switch to the right column. This will now automatically move to the second
%% page if the content is too long.
\switchcolumn

%\cvsection{My Life Philosophy}
%
%\begin{quote}
%``Something smart or heartfelt, preferably in one sentence.''
%\end{quote}

%\cvsection{Most Proud of}
%
%\cvachievement{\faTrophy}{Fantastic Achievement}{and some details about it}
%
%\divider
%
%\cvachievement{\faHeartbeat}{Another achievement}{more details about it of course}
%
%\divider
%
%\cvachievement{\faHeartbeat}{Another achievement}{more details about it of course}

\cvsection{Education}

\cvevent{Master's in Computer Science and Cybersecurity}{University of Lille}{September 2023  -- August 2025}{}

\divider

\cvevent{Computer Science Training}{École 42}{September 2013  -- June 2014}{}

\divider

\cvevent{Bachelor's in Computer Science}{Pierre and Marie Curie University (Paris VI, now Sorbonne University)}{September 2009  -- June 2013}{}

\divider

\cvevent{Self-taught Learning}{127.0.0.1}{December 1991 -- Present}{}

Site du Zéro, Developpez.net, Coursera, IRC, Stack Overflow, ArchWiki, /usr/bin/man

\medskip

\cvsection{Skills}

\cvtag{GNU/Linux}
\cvtag{TCP/IP}
\cvtag{HTML}
\cvtag{CSS}
\cvtag{JS}
\cvtag{TS}
\cvtag{React}
\cvtag{Vue}
\cvtag{Angular}
\cvtag{Elm}
\cvtag{Backbone.js}
\cvtag{Express.js}
\cvtag{SQL (PostgreSQL, MariaDB)}

\cvtag{Bash}
\cvtag{Ruby (+ Ruby on rails, sinatra)}
\cvtag{Python (+ Django and Flask)}
\cvtag{PHP (+ Laravel)}
\cvtag{.NET}
\cvtag{C\#}

\cvtag{Elixir (+ Phoenix)}
\cvtag{Java (+ Spring)}
\cvtag{OCaml}
\cvtag{ReScript}
\cvtag{Perl}
\cvtag{Java}
\cvtag{Lua}
\cvtag{C}
\cvtag{C++}
\cvtag{Rust}
\cvtag{Nix}
\cvtag{Docker}
\cvtag{Ansible}
%\divider\smallskip
\cvtag{Git}
\cvtag{Vim}
%\cvtag{Emacs}

%\divider\smallskip
%
%\cvtag{C++}
%\cvtag{Embedded Systems}\\
%\cvtag{Statistical Analysis}

\cvsection{Languages}

\cvskill{French}{5}
\divider

\cvskill{English}{4.5}
\divider

\cvskill{Spanish}{3}
\divider

\medskip


\end{paracol}

\end{document}
